\chapter{Stephen Smale}\label{ch1966stephensmale:stephen-smale-fields-medal-1966-a-comprehensive-report}

\begin{center}\rule{0.5\linewidth}{0.5pt}\end{center}

\textbf{Prepared:} July 2026
\textbf{Subject:} Stephen Smale (b. July 15, 1930, Flint, Michigan, USA)
\textbf{Position:} Professor, University of California, Berkeley
\textbf{Doctoral Advisor:} Raoul Bott (University of Michigan)
\textbf{Fields Medal Citation:} \emph{``Worked on differential topology and the theory of dynamical systems; contributed important advances to the theory of the classification of differentiable manifolds higher dimensional and the problem of the stability of differentiable mappings. Proved Poincaré's conjecture for dimensions $n \geq 5$.''}

\begin{center}\rule{0.5\linewidth}{0.5pt}\end{center}

\begin{prereqs}
Differential topology (smooth manifolds, transversality, Morse theory), basic algebraic topology (homotopy groups, homology, fundamental group), smooth dynamical systems (flows, fixed points, invariant sets).
\end{prereqs}

\begin{objectives}
Prove the $h$-cobordism theorem using Morse theory and the Whitney trick; derive the generalized Poincaré conjecture for $n \geq 5$ \cite{smale1961}; construct the Smale horseshoe map and characterize chaotic dynamics through symbolic coding.
\end{objectives}

\section{Problem History}

The Poincaré conjecture, formulated by Henri Poincaré in 1904, asks whether every simply connected closed 3-manifold is homeomorphic to the 3-sphere $S^3$.

Poincaré's original formulation was imprecise, and subsequent clarifications only underscored the problem's difficulty.

For over half a century, the conjecture stood as one of the most famous unsolved problems in all of mathematics.

It became the touchstone of algebraic topology, and progress on it was measured in decades.

In dimensions one and two, the classification of surfaces was well understood through the work of Möbius, Riemann, and others.

In dimension three, the conjecture resisted all attempts.

However, in higher dimensions, the existence of extra topological room---formalized by the Whitney trick---suggested that higher-dimensional cases might be more tractable.

This observation, combined with new tools in differential topology, provided the key insight for Smale's breakthrough.

The classification of manifolds sought to determine when two smooth manifolds are diffeomorphic.

Milnor's 1956 discovery of exotic smooth structures on $S^7$ demonstrated that a single topological manifold can carry multiple inequivalent smooth structures.

The h-cobordism theorem, if proved, would classify simply connected manifolds in high dimensions up to diffeomorphism, providing a complete structural theory.

In dynamical systems, the dominant paradigm was that structural stability was generic.

Poincaré had discovered sensitive dependence in the restricted three-body problem, but this was considered exceptional.

Smale's work on the horseshoe revealed that chaotic behavior is ubiquitous and mathematically structured, overturning decades of received wisdom.

\section{How Was It Solved}

The \textbf{h-cobordism theorem} (Smale, 1962): let $W$ be a compact, simply connected $(n+1)$-manifold with boundary $\partial W = M_0 \sqcup M_1$, where $n \geq 5$.

If the inclusions $M_i \hookrightarrow W$ are homotopy equivalences, then $W \cong M_0 \times [0,1]$.

The theorem is the precise higher-dimensional analogue of the statement that a cylinder is trivial.

The proof uses Morse theory.

Choose a Morse function $f: W \to [0,1]$ with $f^{-1}(i) = M_i$.

Since $W$ is an h-cobordism, the Morse inequalities force the critical points to satisfy cancellation relations.

The key step is canceling critical points in adjacent indices using the \textbf{Whitney trick}: in dimension $n \geq 5$, two embedded disks of dimensions $p$ and $q$ with $p + q < n$ can be pushed apart by a Whitney disk.

This allows all critical points to be canceled, leaving $f$ without critical points in the interior, so $W \cong M_0 \times [0,1]$.

The \textbf{generalized Poincaré conjecture} for $n \geq 5$ follows immediately: if $\Sigma^n$ is a homotopy $n$-sphere, removing two open balls yields an h-cobordism $W$ between $S^{n-1}$ and $S^{n-1}$.

By the h-cobordism theorem, $W \cong S^{n-1} \times [0,1]$, so $\Sigma^n \cong S^n$.

This settles the Poincaré conjecture in all dimensions except $n = 3$ and $n = 4$.

The \textbf{Smale horseshoe} is a map $f: D \to \mathbb{R}^2$ on a square $D$.

It contracts vertically by a factor $\delta < 1/2$, stretches horizontally by a factor $\lambda > 2$, and folds the result back into $D$ in the shape of a horseshoe.

The invariant set $\Lambda = \bigcap_{n \in \mathbb{Z}} f^n(D)$ is a Cantor set, and $f|_\Lambda$ is conjugate to the shift $\sigma: \Sigma_2 \to \Sigma_2$ on two symbols.

The shift $\sigma(s_0 s_1 s_2 \cdots) = s_1 s_2 s_3 \cdots$ exhibits dense orbits, dense periodic orbits, sensitive dependence on initial conditions, and topological entropy $h_{\text{top}}(f|_\Lambda) = \log 2$.

The construction of the horseshoe is explicit: given any point $x \in \Lambda$, its itinerary under $f$ determines a unique infinite sequence in $\Sigma_2$, and the conjugacy maps $x$ to this sequence.

The horseshoe is structurally stable: any map sufficiently $C^1$-close to $f$ has a conjugate invariant set.

This robustness is what makes the horseshoe the fundamental model of chaotic dynamics.

\section{Impact of This Discovery}

The h-cobordism theorem is the cornerstone of high-dimensional manifold topology.

Together with Kervaire--Milnor and Wall's surgery theory, it classifies simply connected smooth manifolds of dimension $\geq 5$: two such manifolds are diffeomorphic if and only if they share homotopy type and Pontryagin classes.

The three remaining cases of the Poincaré conjecture---$n = 4$ (Freedman, 1982), $n = 3$ (Perelman, 2003), and the smooth $S^4$ conjecture (open)---each required fundamentally different techniques, underscoring the special nature of low dimensions.

The Smale horseshoe founded the field of symbolic dynamics and established the paradigm of deterministic chaos.

The Smale--Birkhoff homoclinic theorem shows that transverse homoclinic points imply horseshoes, and homoclinic intersections occur generically in non-integrable Hamiltonian systems.

Smale's Axiom A and the $\Omega$-stability theorem provided the first rigorous framework for structural stability.

The horseshoe became the prototypical example in every textbook on dynamical systems.

The Morse--Smale complex connects critical point theory to manifold topology, providing a combinatorial approach to homology that influenced Floer homology and symplectic geometry.

Smale's sphere eversion paradox---turning $S^2$ inside out in $\mathbb{R}^3$ with self-intersections---demonstrated the power of the h-principle and geometric topology.

\section{Modern Views and Improvements}

The \textbf{s-cobordism theorem} extends the h-cobordism theorem to non-simply connected manifolds: for $n \geq 5$, an h-cobordism is trivial iff the Whitehead torsion vanishes in the Whitehead group $\text{Wh}(\pi_1(M_0))$.

Surgery theory (Browder, Novikov, Wall) classifies all compact manifolds via the surgery exact sequence, which reduces the classification to algebraic computations in L-groups and K-theory.

In dimension four, Freedman proved the topological Poincaré conjecture using Casson handles, while Donaldson showed smooth 4-manifolds are far more rigid.

The smooth $S^4$ conjecture---whether $S^4$ admits exotic smooth structures---remains one of the deepest open questions in topology.

The exotic smooth structures on $\mathbb{R}^4$, discovered by Freedman and Donaldson, are unique among all dimensions.

Hyperbolic dynamics (Bowen, Ruelle) and non-uniformly hyperbolic dynamics (Pesin, Katok) extend the horseshoe paradigm to systems that are hyperbolic only on average.

The Palis conjecture asks whether generic smooth systems are structurally stable or have finitely many attractors.

Smooth ergodic theory builds directly on Smale's foundational insights.

Topological data analysis uses Morse theory and persistent homology to extract features from data, connecting Smale's insights to modern computational geometry.

The theory of strange attractors, including the Lorenz and Hénon attractors, builds on the horseshoe as the fundamental model of chaotic behavior.

Smale's later work on the complexity of algorithms and neural computation extended his influence far beyond topology.

\begin{center}\rule{0.5\linewidth}{0.5pt}\end{center}

\section*{Summary}
\addcontentsline{toc}{section}{Summary}

Smale proved the h-cobordism theorem and the generalized Poincaré conjecture for $n \geq 5$, establishing the foundations of high-dimensional manifold topology.

His discovery of the Smale horseshoe and his work on Axiom A dynamics revealed the mathematical structure of chaotic systems, founding the modern theory of dynamical systems.
